\documentclass[11pt,a4paper]{article}
\usepackage[T1]{fontenc}
\usepackage[utf8]{inputenc}
\usepackage{lmodern,amsmath,amssymb,amsthm,mathtools,bm}
\usepackage[a4paper,margin=27mm,headheight=14pt]{geometry}
\usepackage{microtype,booktabs,longtable,array,enumitem,xcolor}
\usepackage{fancyhdr,hyperref}
\definecolor{ink}{HTML}{18374A}
\definecolor{accent}{HTML}{236B80}
\hypersetup{colorlinks=true,linkcolor=ink,citecolor=accent,urlcolor=accent,
 pdftitle={Complex Transseries Reversion at q-Cusps},
 pdfauthor={Research article prepared for Vladimir Reshetnikov}}
\pagestyle{fancy}\fancyhf{}
\fancyhead[L]{\small Complex transseries reversion at $q$-cusps}
\fancyhead[R]{\small October 2026}\fancyfoot[C]{\thepage}
\setlength{\parskip}{4pt}\setlength{\parindent}{1.2em}
\setlist{itemsep=4pt,topsep=5pt}
\numberwithin{equation}{section}
\newtheorem{theorem}{Theorem}[section]
\newtheorem{proposition}[theorem]{Proposition}
\newtheorem{lemma}[theorem]{Lemma}
\newtheorem{corollary}[theorem]{Corollary}
\theoremstyle{definition}\newtheorem{definition}[theorem]{Definition}
\newtheorem{example}[theorem]{Example}
\theoremstyle{remark}\newtheorem{remark}[theorem]{Remark}
\newcommand{\C}{\mathbb C}\newcommand{\Q}{\mathbb Q}\newcommand{\Z}{\mathbb Z}
\newcommand{\N}{\mathbb N}\newcommand{\R}{\mathbb R}
\newcommand{\ii}{\mathrm i}\newcommand{\e}{\mathrm e}
\newcommand{\Log}{\operatorname{Log}}\newcommand{\Li}{\operatorname{Li}}
\newcommand{\ord}{\operatorname{ord}}\newcommand{\id}{\operatorname{id}}
\newcommand{\dd}{\,\mathrm d}\newcommand{\eps}{\varepsilon}
\newcommand{\Poch}{\mathcal P}\newcommand{\F}{\mathcal F}
\newcommand{\calO}{\mathcal O}\newcommand{\supp}{\operatorname{supp}}
\newcommand{\abs}[1]{\left|#1\right|}
\newcommand{\coef}[2]{[#1^{#2}]}
\newcommand{\vct}[1]{\boldsymbol{#1}}
\newcommand{\repo}{\texttt{VladimirReshetnikov/ProveIt}}
\begin{document}
\hypersetup{pageanchor=false}
\begin{titlepage}
\vspace*{13mm}
{\color{accent}\large RESEARCH ARTICLE \hfill 4 OCTOBER 2026}\par
\vspace{16mm}
{\color{ink}\Huge\bfseries Complex Transseries\\[4pt] Reversion at $q$-Cusps}\par
\vspace{8mm}
{\Large Convergent normal forms, infinitely flat limits,\\[3pt]
logarithmic--Puiseux inverses, and recovery from cusp data}\par
\vspace{12mm}\hrule\vspace{7mm}
{\large Prepared for Vladimir Reshetnikov}\par
\vspace{9mm}
\begin{abstract}
We develop a branch-resolved inversion calculus for complex transseries and
apply it to $q$-products at rational boundary points. The analytic core is a
convergent, multivariable implicit-function construction for
$Y=Ct^\beta\e^{-A/t}U(t,\e^{-B_1/t},\ldots,\e^{-B_d/t})$.
It retains complete exponential sectors, gives explicit coefficient
extraction and residual bounds, and converts ratios of actions into powers
of the inverse variable. A second scaling handles power and regular cores
whose exponentially small perturbations have large logarithmic derivatives.

For every finite Euler-product quotient
$\prod_{\delta\mid N}(q^\delta;q^\delta)_\infty^{r_\delta}$,
we derive an exact analytic normal form at every root of unity. Three
explicit linear invariants determine the inverse regime. When all three
vanish, the quotient has a nonconstant but infinitely flat finite limit;
its inverse is described by a convergent Puiseux series inside a logarithmic
denominator. We prove a four-factor minimum at $q=1$ and realize every finite correction
ramification degree with four factors. We compute two examples and recover every exponent $r_\delta$ from the leading
exponential data at divisor cusps by a double M\"obius inversion. Additional
results treat nonmodular shifted factorials, finite Gaussian multinomials,
and the transport of sectorial jumps through inversion.

All stated local theorems have conventional proofs. Classical modular,
Lagrange, and resurgent-closure results are attributed; the proposed
contribution is their explicit branch-aware synthesis and the accompanying
classification and reconstruction formulas. Worldwide novelty is not
established. The included computations are reproducible checks, not
proof-assistant or interval certificates.
\end{abstract}
\vfill
{\small\color{ink}Self-contained source and bibliography.\par Repository snapshot:\par
\texttt{8e9cd6f00e0ac79c4425ff6abdfa64b497cf3f41}.}
\end{titlepage}
\hypersetup{pageanchor=true}\setcounter{page}{1}
\begingroup\small\setlength{\parskip}{0pt}
\tableofcontents
\endgroup\newpage

\section{Scope, source boundary, and the questions answered}
\label{sec:scope}

\subsection{What ``general complex transseries'' can responsibly mean}
Series reversal means compositional inversion, not reversal of the order of
coefficients. Over the complex plane, the choice of sector, logarithm,
ramification, and analytic realization is part of the inversion problem.
There cannot be a single branch-free inverse calculus even for
$Y=\e^{-1/t}$: its inverses are
$-1/(\Log Y+2\pi\ii k)$, $k\in\Z$.

We distinguish three levels throughout. A \emph{formal} statement takes
place in a specified complete filtered algebra. A \emph{sectorial}
statement concerns actual holomorphic functions with uniform estimates on
smaller sectors. A \emph{global continuation} statement requires transition
maps and their domains. No implication from the first level to the third
is assumed. In particular, attaching complex coefficients to a real
transseries field does not by itself define all oscillatory compositions
or all analytic sums.

The generality proved here is a preparation principle: after a dominant
phase is inverted and the remaining small quantities are promoted to
independent parameters, a substantial class of complex transseries becomes
a convergent analytic implicit-function problem. The formal version also
works in other complete filtered algebras when its contraction hypotheses
are verified. We do not assert that every conceivable Hahn support,
accumulating action set, or divergent transseries admits such a preparation.

\subsection{Relation to ProveIt and to existing literature}
The inspected repository snapshot is the commit printed on the title page.
The canonical transseries directory records a large inversion volume,
a combinatorial companion, and separate research arrivals; it explicitly
distinguishes human arguments from exact Lean counterparts
\cite{repo}. The companion already contains phase-coordinate reversal,
analytic exponential-tail inversion, and applications to $q$-products.
Consequently, neither ``use Lagrange inversion'' nor ``invert an exponential
tail'' is presented here as a new foundational discovery.

The inspected \emph{Certified Inversion Through the $q\to1$ Transition}
README describes signed remainders, Stieltjes--Pad\'e bounds, optimal
truncation, and a cancellation window for Gaussian products
\cite{repoq}. The \emph{Beyond Finite-Action Folds} README concerns
countable-action critical models and Hahn inverses \cite{repohahn}.
Those are different problems from classifying an Euler quotient's finite
flat limit at every rational cusp. This is a targeted crosswalk based on
the listed documentation, not a claim to have audited all pages of the
canonical volume or every incoming article.

Formal and analytic nonlinear operations on resurgent functions, including
functional and implicit inversion, are established subjects. Sauzin proves
precise closure and summability theorems in an appropriate resurgent
setting \cite{sauzin}. We use those results as background, not as an open
conjecture supposedly solved here. The modular transformation of Dedekind's
eta function is classical \cite{dlmf}. All-order $q$-factorial asymptotics
involving polylogarithms and Bernoulli polynomials occur in McIntosh's work
\cite{mcintosh}, and root-of-unity asymptotics are central to the work of
Garoufalidis and Zagier \cite{gz}.

Two recent sources materially sharpen the source boundary. Fantini and
Rella analyze resurgence and summability of $q$-Pochhammer symbols and
weighted combinations \cite{fr}. Arabi Ardehali and Rosengren give a
gamma-product representation and asymptotic applications, with a version
of record published on 19 September 2026 \cite{ar}. We do not claim a new
$q$-Pochhammer product formula or a solution of their general summability
questions. Our exact results use the modular Euler-product subclass;
Section~\ref{sec:shifted} separately treats an asymptotic, generally
nonconvergent, shifted-factorial example.

\subsection{A precise result ledger}
The following concrete questions are answered in this article.
\begin{enumerate}[label=\textbf{Q\arabic*.},leftmargin=14mm]
\item Can inversion preserve full exponential sectors rather than just
algebraic coefficients? Theorem~\ref{thm:prepared} gives a convergent
multivariable lift, sector-preserving truncation, and action transport.
\item What replaces ordinary Puiseux inversion at a nonconstant finite
limit invisible to every power of $t$? Theorem~\ref{thm:flat} gives all
local sheets, their monodromy, and an explicit error law.
\item Which finite Euler quotients have such a limit at a specified root
of unity? Theorem~\ref{thm:classification} gives a necessary-and-sufficient
three-invariant test and an inverse in each complementary regime.
\item How few factors can produce flatness at $q=1$? Four distinct
scales are necessary and sufficient, with an explicit kernel formula
(Theorem~\ref{thm:four}).
\item Can all rational cusps conceal the quotient's exponent vector?
No: Theorem~\ref{thm:recover} gives exact recovery from divisor-cusp
exponential data.
\item Is finite correction ramification bounded by the number of product
factors? No: four factors realize every positive degree, with an exact
intrinsic-degree criterion (Theorem~\ref{thm:ramification}).
\end{enumerate}
These are explicitly posed local questions. Answering them is not a claim
of resolving an independently documented major open problem. The proofs
and formulations are research contributions of this draft; publication
priority and independent validation remain separate matters.

\section{Sectorial preparation and reliable inverse certificates}
\label{sec:foundations}

\subsection{Small actions and admissible sectors}
A typical asymptotic sector is
\[
 S(\theta,\alpha,r)=\{0<|t|<r:\ |\arg t-\theta|<\alpha\},
\]
with its argument understood on a chosen logarithmic cover. Bounds are
uniform on closed, strictly smaller subsectors. A quantity
$\e^{-B/t}$ is small only where $\Re(B/t)>0$; it is uniformly exponentially
small where $\Re(B/t)\ge c/|t|$ for some $c>0$.

For finitely many actions $\lambda_j$, assume a sector of large $x$ on
which $\Re(\lambda_jx)\ge c|x|$ for every $j$. Their additive semigroup is
locally finite when measured by this positive real projection: only
finitely many multi-indices have projection at most a given bound.
This permits resonant actions to be collected without infinite
coefficient multiplicities. It does not supply a total asymptotic ordering
on every direction. In particular, a dominance wall is not automatically
a Stokes jump.

\begin{lemma}[Complete filtered contraction]\label{lem:formal}
Let $\mathcal A$ be a complete separated $\C$-algebra with descending
multiplicative ideals $I^n$, $n\ge0$. Suppose a map
$T:I\to I$ is defined by admissible algebraic or Taylor operations and
\[
 u-v\in I^n\quad\Longrightarrow\quad T(u)-T(v)\in I^{n+1}.
\]
Then $u=T(u)$ has a unique solution in $I$. Its residue modulo $I^{n+1}$
is obtained after $n$ Picard steps from $0$.
\end{lemma}
\begin{proof}
Put $u_0=0$ and $u_{n+1}=T(u_n)$. Induction gives
$u_{n+1}-u_n\in I^{n+1}$. Completeness gives a limit $u$; the displayed
condition implies continuity and hence $u=T(u)$. For two fixed points,
$u-v\in I$ improves successively to every $I^n$, so separatedness gives
$u=v$. The same induction gives the residue assertion.
\end{proof}

The admissibility clause is essential. One must not differentiate,
compose, or sum arbitrary Hahn families merely because the resulting
expression is suggestive. For analytic parameter germs we use
$\C[[v,Q_1,\ldots,Q_d]]$ and its maximal-ideal filtration; all operations
below are then explicit and admissible.

\subsection{A complex residual certificate with existence}
\begin{lemma}[Derivative-close-to-one certificate]\label{lem:certificate}
Let $K$ be holomorphic on a neighborhood of the closed disk
$\overline D(z_*,r)$. Suppose
\[
 \sup_{\overline D(z_*,r)}|K'(z)-1|\le\kappa<1,
 \qquad |K(z_*)|<(1-\kappa)r.
\]
Then $K$ has a unique zero $z$ in that disk and
\begin{equation}\label{eq:certificate}
 |z-z_*|\le \frac{|K(z_*)|}{1-\kappa}.
\end{equation}
\end{lemma}
\begin{proof}
The map $T(w)=w-K(w)$ has Lipschitz constant at most $\kappa$ on the
convex disk. Moreover,
$|T(w)-z_*|\le\kappa|w-z_*|+|K(z_*)|<r$.
The contraction theorem gives a unique fixed point. At that point,
$|z-z_*|\le\kappa|z-z_*|+|K(z_*)|$, proving the bound.
\end{proof}
A lower bound on $|K'|$ alone is not a substitute for this argument in the
complex plane: it does not prove injectivity on an arbitrary region.
The lemma supplies both existence and the backward-to-forward error
conversion needed by an inverse algorithm.

\section{A convergent lift for prepared exponential transseries}
\label{sec:prepared}

\subsection{The exact normal form}
Fix $A,C\ne0$, $\beta\in\C$, and actions $B_1,\ldots,B_d\in\C$.
Let $U(t,\vct E)$ be holomorphic near $(0,\vct0)$ with $U(0,\vct0)=1$.
Consider
\begin{equation}\label{eq:prepared-map}
 Y=Ct^\beta\e^{-A/t}
 U(t,\e^{-B_1/t},\ldots,\e^{-B_d/t}).
\end{equation}
Choose all logarithms coherently so that, with $x=A/t$,
$\Log t=\Log A-\Log x$. Put
\begin{equation}\label{eq:core}
 L=\Log\frac{CA^\beta}{Y},\qquad
 x_0+\beta\Log x_0=L,\qquad
 v=x_0^{-1},\quad \lambda_j=B_j/A,\quad Q_j=\e^{-\lambda_jx_0}.
\end{equation}
On a closed large-$L$ sector away from the logarithmic cut, the root with
$x_0\sim L$ exists and is unique in an $O(\log|L|)$ disk around $L$.
Indeed, $x\mapsto L-\beta\Log x$ maps such a disk to itself and has
Lipschitz constant $O(|L|^{-1})$. Thus
\begin{equation}\label{eq:x0-expansion}
 x_0=L-\beta\Log L+O(\log|L|/|L|).
\end{equation}
When $\beta\ne0$, the identity
$x_0=\beta W(\e^{L/\beta}/\beta)$ is available, but the branch is the one
selected by \eqref{eq:core}; it need not be the principal $W$ branch.

\begin{theorem}[Analytic parameter lift]\label{thm:prepared}
There is a unique holomorphic germ $D(v,\vct Q)$ at the origin, with
$D(0,\vct0)=0$, satisfying
\begin{align}
 D={}&-\beta\log(1+vD)\notag\\
 &+\log U\left(\frac{Av}{1+vD},
 Q_1\e^{-\lambda_1D},\ldots,Q_d\e^{-\lambda_dD}\right).
 \label{eq:D}
\end{align}
In every sector in which $|x_0|\to\infty$ and
$\Re(\lambda_jx_0)\ge c|x_0|$, the inverse of
\eqref{eq:prepared-map} near $A/x_0$ is exactly
\begin{equation}\label{eq:prepared-inverse}
 t=\frac{A}{x_0+D(x_0^{-1},\e^{-\lambda_1x_0},\ldots,
 \e^{-\lambda_dx_0})}.
\end{equation}
Its multivariable Taylor expansion is convergent, and no exponential
action outside the additive semigroup generated by the $\lambda_j$
is needed in the $x_0$ coordinate.
\end{theorem}
\begin{proof}
Subtract the logarithm of \eqref{eq:prepared-map} from
\eqref{eq:core} and write $x=x_0+D$. The resulting equation is
\eqref{eq:D}. Its left side minus right side is holomorphic in
$(D,v,\vct Q)$ near zero. Its derivative with respect to $D$ there is
$1$: the logarithmic term has a factor $v$, the first argument of $U$
has a factor $v$, and each other argument has a factor $Q_j$.
The analytic implicit-function theorem gives the germ and uniqueness.
The right side also has positive parameter order and raises the order of
a difference in $D$, so Lemma~\ref{lem:formal} gives the same formal germ.

Under the sector hypotheses the actual parameters tend to the origin.
Substitution into the convergent germ solves the original equation with
its chosen logarithms. The prepared equation has derivative close to one
on a fixed small $D$ disk, giving local uniqueness by
Lemma~\ref{lem:certificate}. Finally, a Taylor monomial
$v^k\vct Q^{\vct n}$ evaluates to
$x_0^{-k}\e^{-(\vct n\cdot\vct\lambda)x_0}$.
Convergence and the positive projection condition justify grouping
resonant actions. Analytic reciprocation in \eqref{eq:prepared-inverse}
preserves the same action semigroup.
\end{proof}

\subsection{Coefficient extraction and sector-preserving truncation}
Write the right side of \eqref{eq:D}, with $D$ replaced by $z$, as
$\Phi(z;v,\vct Q)$. Classical Lagrange inversion with an auxiliary
parameter gives
\begin{equation}\label{eq:lagrange-D}
 D=\sum_{n\ge1}\frac1{n!}
 \left.\frac{\partial^{n-1}}{\partial z^{n-1}}
 \Phi(z;v,\vct Q)^n\right|_{z=0}.
\end{equation}
This is a coefficientwise identity in the parameter filtration: every
summand has order at least $n$. It is also an analytic identity for
sufficiently small parameters. To see the latter directly, on a fixed
small $z$ disk the whole $\Phi$ is uniformly small, so the auxiliary
parameter can vary in a disk of radius larger than $1$.

Suppose $D$ is bounded by $M$ on a closed polydisk of radius $R$, slightly
inside its domain. If $D_{\le n}$ denotes total-parameter-degree
truncation and
$s=(|v|+\sum_j|Q_j|)/R<1$, Cauchy's estimate yields
\begin{equation}\label{eq:total-tail}
 |D-D_{\le n}|\le M\frac{s^{n+1}}{1-s}.
\end{equation}
The sum of monomials of degree $k$ is bounded by the $k$th power of the
sum of their absolute variables; no missing multinomial factor is needed.

For exponential accuracy, a different grouping is preferable:
\begin{equation}\label{eq:action-block}
 D(v,\vct Q)=\sum_{\vct n\in\N^d}D_{\vct n}(v)\vct Q^{\vct n}.
\end{equation}
Keep the coefficient functions $D_{\vct n}(v)$ intact. Uniformly for $v$
in a smaller disk, the tail after $|\vct n|\le n$ is bounded by the
analogue of \eqref{eq:total-tail} with
$s=\sum_j|Q_j|/R$. Truncating every coefficient function to a fixed
algebraic order instead leaves an error much larger than an exponential
sector. Thus a finite algebraic approximation does not automatically
certify a beyond-all-orders claim.

For an approximate displacement $\widetilde D$, apply
Lemma~\ref{lem:certificate} to the left side minus right side of
\eqref{eq:D}. Once both displacements are in its disk,
\begin{equation}\label{eq:t-residual}
 |t-\widetilde t|
 \le\frac{|A|\,|K(\widetilde D)|}
 {(1-\kappa)|x_0+D|\,|x_0+\widetilde D|}.
\end{equation}
This is an explicit certificate, including a root-existence test.

\subsection{Actions become inverse-variable exponents}
The core identity gives the exact branch-dependent relation
\begin{equation}\label{eq:action-transport}
 Q_j=\e^{-\lambda_jL}x_0^{\beta\lambda_j}
 =\left(\frac{Y}{CA^\beta}\right)^{\lambda_j}
 x_0^{\beta\lambda_j}.
\end{equation}
Thus $B_j/A$, not $B_j$ alone, is the exponent visible in the inverse
variable. Logarithmic dressing comes from the core $x_0$.
All powers in this formula are defined through the same logarithm as $L$.

\begin{proposition}[A sharp one-action ramification test]\label{prop:irrational}
For $a,b>0$, let
$Y=c\e^{-a/t}(1+\e^{-b/t})$, $c\ne0$, and set
$\lambda=b/a$, $L=\Log(c/Y)$. Its inverse is
$t=a/(L+D((Y/c)^\lambda))$, where $D$ is convergent near zero and
\[
 D(Q)=Q-(\lambda+\tfrac12)Q^2+O(Q^3).
\]
The correction to the inverse phase is a convergent Puiseux germ in $Y/c$
after finite ramification if and only if $\lambda\in\Q$.
\end{proposition}
\begin{proof}
The implicit equation is $D=\log(1+Q\e^{-\lambda D})$.
It gives the displayed coefficients and $D'(0)=1$. Rational $\lambda$
clearly suffices. Conversely, a nonzero convergent Puiseux germ has a
rational first exponent. The first exponent here is exactly $\lambda$,
so irrational $\lambda$ is impossible. The logarithm in $L$ is retained;
the assertion concerns its analytic phase correction, not removal of
logarithmic monodromy.
\end{proof}

\section{Flat perturbations of power and regular cores}
\label{sec:power}
When the leading essential action vanishes, Theorem~\ref{thm:prepared}
is no longer the right coordinate chart. Consider instead
\begin{equation}\label{eq:power-map}
 Y=Ct^\beta\e^{Bt}H(\e^{-a/t}),\qquad a>0,\quad H(0)=1,
\end{equation}
with $H$ holomorphic near zero. Assume $(\beta,B)\ne(0,0)$.
Let $\tau$ be a chosen small root of the \emph{exact core equation}
\begin{equation}\label{eq:power-core}
 Y=C\tau^\beta\e^{B\tau}.
\end{equation}
For $\beta\ne0$, choose a branch of
$z=(Y/C)^{1/\beta}$; then
$\tau=\beta W((B/\beta)z)/B$ near zero when $B\ne0$, and
$\tau=z$ when $B=0$. For $\beta=0$, $B\ne0$,
$\tau=B^{-1}\Log(Y/C)$.

\begin{theorem}[Slope-adapted analytic lift]\label{thm:power}
Put $\sigma=1$ if $\beta\ne0$, and $\sigma=2$ if
$\beta=0$, $B\ne0$. There is a holomorphic germ $D(\tau,R)$ with
$D(\tau,0)=0$ for which the exact local inverse is
\begin{equation}\label{eq:power-inverse}
 t=\tau+\tau^2D\left(\tau,\frac{\e^{-a/\tau}}{\tau^\sigma}\right).
\end{equation}
This is a convergent transseries on closed subsectors of
$\Re(1/\tau)>0$. If $H(Q)=1+h_1Q+O(Q^2)$, its first displacement is
\begin{equation}\label{eq:power-first}
 t-\tau\sim
 \begin{cases}
 -\dfrac{h_1}{\beta}\tau\e^{-a/\tau},&\beta\ne0,\\[4pt]
 -\dfrac{h_1}{B}\e^{-a/\tau},&\beta=0,\ B\ne0,
 \end{cases}
\end{equation}
provided $h_1\ne0$.
\end{theorem}
\begin{proof}
Write $t=\tau(1+\tau D)$. Then
\[
 \e^{-a/t}=\e^{-a/\tau}
 \exp\left(\frac{aD}{1+\tau D}\right).
\]
Subtract the logarithm of the core equation. For $\beta\ne0$, divide
by $\tau$ to obtain
\begin{align*}
 0={}&\beta\frac{\log(1+\tau D)}{\tau}+B\tau D\\
 &+\frac1\tau\log H\left(\tau R
 \exp\left(\frac{aD}{1+\tau D}\right)\right).
\end{align*}
All apparent singularities at $\tau=0$ are removable. The derivative
with respect to $D$ at $(D,\tau,R)=0$ is $\beta\ne0$.
For $\beta=0$, divide instead by $\tau^2$:
\[
 0=BD+\frac1{\tau^2}\log H\left(\tau^2R
 \exp\left(\frac{aD}{1+\tau D}\right)\right).
\]
The derivative is now $B\ne0$. The analytic implicit-function theorem
proves both cases; setting $R=0$ gives $D=0$. Exponential decay ensures
$\e^{-a/\tau}/\tau^\sigma\to0$ in the stated subsectors. Linearization
in $R$, followed by $\tau\to0$, gives \eqref{eq:power-first}.
\end{proof}
The division by $\tau$ or $\tau^2$ is not cosmetic. It absorbs the
large logarithmic derivative of the exponential into the small parameter. Expanding
only in $Q$ while pretending that all coefficients stay uniformly bounded
at $\tau=0$ would miss this scaling.

\section{Infinitely flat finite limits and logarithmic--Puiseux inversion}
\label{sec:flat}

\subsection{All sheets and a convergent denominator}
Suppose
\begin{equation}\label{eq:flat-map}
 Y(t)=Y_*+c\e^{-ma/t}U(\e^{-a/t}),\qquad
 a>0,\quad c\ne0,\quad m\in\N_{>0},\quad U(0)=1,
\end{equation}
where $U$ is holomorphic near zero. Let $s$ and $\log s$ be fixed on a
simply connected punctured neighborhood by
\[
 s=\left(\frac{Y-Y_*}{c}\right)^{1/m},\qquad
 \log s=\frac1m\Log\frac{Y-Y_*}{c}.
\]

\begin{theorem}[Logarithmic--Puiseux inverse atlas]\label{thm:flat}
Let $\psi$ be the analytic inverse at zero of
$f(Q)=Q\,U(Q)^{1/m}$, where $U^{1/m}(0)=1$, and let
$V(z)=\log(\psi(z)/z)$ with $V(0)=0$. Then all solutions with
$|\e^{-a/t}|$ sufficiently small are
\begin{equation}\label{eq:flat-inverse}
 t_\ell(Y)=-\frac{a}{\Lambda_\ell+V(s_\ell)},\qquad
 s_\ell=\e^{2\pi\ii\ell/m}s,\quad
 \Lambda_\ell=\log s+\frac{2\pi\ii\ell}{m},\quad \ell\in\Z.
\end{equation}
The coefficients are explicit:
\begin{equation}\label{eq:psi-coeff}
 \psi(z)=\sum_{n\ge1}\frac{z^n}{n}
 [Q^{n-1}]\,U(Q)^{-n/m}.
\end{equation}
A positive circuit of $Y$ around $Y_*$ continues sheet $\ell$ to
sheet $\ell+1$.
\end{theorem}
\begin{proof}
The map $f$ is holomorphic, $f(0)=0$, and $f'(0)=1$. Hence it has a
unique local analytic inverse $\psi$. With $E=\e^{-a/t}$,
\eqref{eq:flat-map} becomes $f(E)^m=s^m$, so
$E=\psi(\e^{2\pi\ii j/m}s)$ for one $j$ modulo $m$.
The logarithms of this nonzero $E$ are
\[
 \log s+2\pi\ii j/m+V(\e^{2\pi\ii j/m}s)+2\pi\ii k,
 \qquad k\in\Z.
\]
Writing $\ell=j+mk$ gives exactly \eqref{eq:flat-inverse}, with no
independent or double-counted branch labels. Conversely, exponentiating
that equation recovers $E$ and then $Y$. Formula~\eqref{eq:psi-coeff}
is ordinary Lagrange inversion. Under a positive circuit of $Y-Y_*$,
$s$ continues to $\e^{2\pi\ii/m}s$ and $\log s$ increases by
$2\pi\ii/m$; substitution proves the monodromy assertion.
\end{proof}

Although the cusp lies on the boundary of the original $q$ disk, the
inverse branches are actual holomorphic functions on slit target
neighborhoods. For small $s$, $\Re(\Lambda_\ell+V(s_\ell))<0$;
therefore $\Re t_\ell>0$. For each fixed $\ell$, $t_\ell\to0$ as
$Y\to Y_*$. The exact formula also permits indices that grow, but claims
about uniform sectorial asymptotics then require a separate bound on the
argument of the denominator.

\subsection{Remainder law and sensitivity}
\begin{corollary}[Explicit inverse remainder]\label{cor:flat-tail}
Suppose $V$ is holomorphic near the closed disk $|z|\le R$,
$|V|\le M$ there with $M>0$, and $V_N$ is its Taylor polynomial through degree $N$.
For $\rho=|s|/R\le1/2$ and $|\Lambda_\ell|\ge2M$,
\begin{equation}\label{eq:flat-tail}
 \left|t_\ell+\frac{a}{\Lambda_\ell+V_N(s_\ell)}\right|
 \le \frac{4aM}{|\Lambda_\ell|^2}
 \frac{\rho^{N+1}}{1-\rho}.
\end{equation}
For fixed sheet, this is
$O(|Y-Y_*|^{(N+1)/m}/|\Log(Y-Y_*)|^2)$.
\end{corollary}
\begin{proof}
Cauchy's bound gives
$|V-V_N|\le M\rho^{N+1}/(1-\rho)$.
Because $V(0)=0$, $|V_N|\le M\sum_{j=1}^N\rho^j\le M$;
also $|V|\le M$. Both denominators have modulus at least
$|\Lambda_\ell|/2$. Subtract the two reciprocals to get
\eqref{eq:flat-tail}.
\end{proof}

Differentiating \eqref{eq:flat-map}, with $E=\e^{-a/t}$, gives
\begin{equation}\label{eq:flat-sensitivity}
 \frac{\dd t}{\dd Y}
 =\frac{t^2}{ma(Y-Y_*)}\bigl(1+O(E)\bigr).
\end{equation}
Thus the inverse is badly conditioned with respect to \emph{absolute}
errors in $Y$ near $Y_*$. In relative displacement units, however,
$(Y-Y_*)t'(Y)/t$ is only of order $1/|\Log(Y-Y_*)|$ on a fixed sheet.
An algorithm should therefore compute $Y-Y_*$ directly, or with enough
precision to survive its cancellation, rather than subtracting two
low-precision numbers close to $Y_*$.

Every map in \eqref{eq:flat-map} has the same algebraic asymptotic
expansion $Y_*+0t+0t^2+\cdots$, irrespective of $a,m,c,U$.
Consequently, algebraic coefficients alone do not determine the inverse,
its leading scale, or its branch structure. The full transseries data do.
This is a failure of information loss under projection to a power series,
not a failure of transseries inversion.

\section{An exact normal form at every rational \texorpdfstring{$q$}{q}-cusp}
\label{sec:cusps}

\subsection{Euler quotients and cusp coordinates}
Write
\[
 \Poch(q)=(q;q)_\infty=\prod_{n\ge1}(1-q^n),\qquad |q|<1.
\]
Let $N\ge1$ and let $r_\delta\in\Z$, $\delta\mid N$, not all vanish.
Our main family is
\begin{equation}\label{eq:euler-quotient}
 \F_{\vct r}(q)=\prod_{\delta\mid N}\Poch(q^\delta)^{r_\delta}.
\end{equation}
Zero exponents may be suppressed. Each factor is nonzero in the unit
 disk, so negative exponents cause no poles there. We call
\eqref{eq:euler-quotient} an \emph{Euler quotient}: converting it to a
literal eta quotient introduces a known elementary power of $q$.
No modular-form congruences on $\vct r$ are imposed.

Choose coprime integers $h,k$, $k>0$, and put
\[
 \zeta=\e^{2\pi\ii h/k},\qquad q=\zeta\e^{-t},\qquad \Re t>0.
\]
The cusp $q=1$ is represented by $(h,k)=(0,1)$. For each active scale set
\begin{align}
 g_\delta&=\gcd(k,\delta),& c_\delta&=k/g_\delta,&
 a_\delta&=\delta h/g_\delta.\label{eq:cusp-integers}
\end{align}
Choose $d_\delta$ with $a_\delta d_\delta\equiv1\pmod{c_\delta}$;
when $c_\delta=1$, take $d_\delta=0$. Define
\begin{equation}\label{eq:dedekind}
 s(d,c)=\sum_{n=1}^{c-1}
 \left(\frac nc-\frac12\right)
 \left(\frac{dn\bmod c}{c}-\frac12\right),\qquad s(d,1)=0.
\end{equation}
For $c>1$ the residues in this sum are nonzero because $d$ is coprime
to $c$. Thus this is the usual Dedekind sum.

The data needed for inversion are
\begin{align}
 R_0&=\sum_{\delta\mid N}r_\delta,&
 R_1&=\sum_{\delta\mid N}\delta r_\delta,&
 \beta&=-R_0/2,& B&=R_1/24,\label{eq:R-data}\\
 A_k&=\frac{\pi^2}{6k^2}
       \sum_{\delta\mid N}\frac{r_\delta g_\delta^2}{\delta},
 & a&=\frac{4\pi^2}{k^2N},&
 m_\delta&=\frac{Ng_\delta^2}{\delta},&
 \rho_\delta&=\e^{-2\pi\ii d_\delta/c_\delta}.\label{eq:A-data}
\end{align}
Every $m_\delta$ is a positive integer. Finally put
\begin{align}
 C_{h,k}
 &=(2\pi)^{R_0/2}
 \prod_{\delta\mid N}(c_\delta\delta)^{-r_\delta/2}
 \exp\left(-\pi\ii\sum_{\delta\mid N}
                 r_\delta s(d_\delta,c_\delta)\right),
 \label{eq:C-data}\\
 H_{h,k}(Q)&=\prod_{\delta\mid N}
                \Poch(\rho_\delta Q^{m_\delta})^{r_\delta}.
 \label{eq:H-data}
\end{align}
The function $H_{h,k}$ is holomorphic and nonzero for $|Q|<1$, and
$H_{h,k}(0)=1$. The number $A_k$ depends only on the cusp denominator;
$C_{h,k}$ and the phases in $H_{h,k}$ also depend on the numerator.

\begin{theorem}[Exact cusp preparation]\label{thm:cusp}
For every $t$ with $\Re t>0$,
\begin{equation}\label{eq:cusp-normal}
 \boxed{\quad
 \F_{\vct r}(\zeta\e^{-t})
 =C_{h,k}\,t^\beta\exp(-A_k/t+Bt)
           H_{h,k}(\e^{-a/t}).\quad}
\end{equation}
Here $t^\beta$ uses the logarithm on the right half-plane that is real
on the positive axis. This is an identity of holomorphic functions, not
merely an asymptotic expansion.
\end{theorem}
\begin{proof}
We record the reduction to the classical eta transformation
\cite[\S27.14]{dlmf} to fix all phases. For one scale omit the subscript
$\delta$ on $a_\delta,c_\delta,d_\delta$, and choose an integer $b$ with
$ad-bc=1$. Set
\[
 z=\frac ac+\frac{\ii\delta t}{2\pi},\qquad
 w=\frac{dz-b}{-cz+a}=-\frac dc+\frac{2\pi\ii}{c^2\delta t}.
\]
Both $z,w$ lie in the upper half-plane. In a compatible square-root
convention the eta law is
\[
 \eta\left(\frac{aw+b}{cw+d}\right)
 =\exp\!\left(\pi\ii\left(\frac{a+d}{12c}-s(d,c)\right)\right)
 [-\ii(cw+d)]^{1/2}\eta(w).
\]
Use $\Poch(\e^{2\pi\ii z})=\e^{-\pi\ii z/12}\eta(z)$ and
$cw+d=2\pi\ii/(c\delta t)$. The result is
\begin{align}
 \Poch(\zeta^\delta\e^{-\delta t})
 ={}&\e^{-\pi\ii s(d,c)}
 \left(\frac{2\pi}{c\delta t}\right)^{1/2}
 \exp\left(-\frac{\pi^2}{6c^2\delta t}+\frac{\delta t}{24}\right)
 \notag\\
 &\hspace{15mm}\cdot
 \Poch\left(\e^{-2\pi\ii d/c}
                   \e^{-4\pi^2/(c^2\delta t)}\right).
 \label{eq:single-cusp}
\end{align}
In particular, the terms $(a+d)/(12c)$ cancel exactly; discarding them
before transforming would give the wrong phase. The square root in
\eqref{eq:single-cusp} is positive for positive $t$ before its displayed
Dedekind multiplier. Multiplying the individual formulas to powers
$r_\delta$ gives \eqref{eq:cusp-normal} and all its constants. The
formula for $c=1$ includes integer translations of $z$ and has the same
cancellation. No modular-form transformation property for the assembled
quotient was needed.
\end{proof}

\subsection{Exact coefficients and termination of the first-action search}
Let $\sigma_1(n)=\sum_{d\mid n}d$. Writing
$\log H_{h,k}(Q)=\sum_{n\ge1}\ell_nQ^n$, direct expansion of
$\log\Poch(z)$ gives
\begin{equation}\label{eq:logH}
 \ell_n=-\frac1n
 \sum_{\substack{\delta\mid N\\m_\delta\mid n}}
 r_\delta m_\delta\,
 \sigma_1(n/m_\delta)\rho_\delta^{\,n/m_\delta}.
\end{equation}
If $H_{h,k}(Q)=\sum_{n\ge0}h_nQ^n$, then
\begin{equation}\label{eq:H-rec}
 h_0=1,\qquad
 nh_n=\sum_{j=1}^n j\ell_jh_{n-j}\quad(n\ge1).
\end{equation}
Thus every coefficient is in an explicitly specified cyclotomic field.
At $q=1$, all phases are $1$ and the coefficients are integers: this
also follows directly from integer powers of the Euler factors.

\begin{lemma}[The transformed tail is nonconstant]\label{lem:H-nonconstant}
For a nonzero finite exponent vector $\vct r$, the function $H_{h,k}$
is not identically $1$. Consequently the search for the first nonzero
coefficient of $H_{h,k}-1$ terminates at every cusp.
\end{lemma}
\begin{proof}
Suppose $H_{h,k}=1$. By analytic continuation in the unit disk it is
identically $1$ throughout its domain. Formula~\eqref{eq:cusp-normal}
then holds without a tail for every real $t>0$. As $t\to+\infty$,
$\F_{\vct r}(\zeta\e^{-t})=1+O(\e^{-\delta_0t})$, where
$\delta_0$ is the smallest active scale. Since $B,\beta,A_k$ are real,
comparison with $C_{h,k}t^\beta\e^{-A_k/t+Bt}$ first gives
$B=0$, then $\beta=0$ and $C_{h,k}=1$, and finally $A_k=0$ by
comparison of the $1/t$ terms. It follows that $\F_{\vct r}=1$.
But its Taylor coefficient at $q^{\delta_0}$ is $-r_{\delta_0}\ne0$;
larger scales cannot contribute to that coefficient. This contradiction
proves the claim.
\end{proof}
The proof is qualitative: it does not give a useful uniform bound for
the first nonzero index at arbitrary cusps. At $q=1$ there is a simple
bound and an exact answer. If $\delta_{\max}$ is the largest active
scale, the smallest $m_\delta=N/\delta$ is unique, so
\begin{equation}\label{eq:first-q1}
 m=\ord_0(H_{0,1}-1)=N/\delta_{\max},\qquad h_m=-r_{\delta_{\max}}.
\end{equation}
Changing the auxiliary common multiple $N$ merely changes the coordinate
$Q$; the physical first action $ma=4\pi^2/\delta_{\max}$ is unchanged.
An apparent Puiseux order can contain redundant ramification if $H$ is
a function of a proper power of $Q$. The complete sheet formula remains
valid; a smaller primitive action coordinate removes that redundancy.

\section{Classification of cusp inverses and a four-factor minimum}
\label{sec:classification}

\begin{theorem}[Four inverse regimes]\label{thm:classification}
For a nontrivial Euler quotient \eqref{eq:euler-quotient} at a fixed cusp,
exactly one of the following applies, with the data of
Theorem~\ref{thm:cusp}.
\begin{enumerate}[label=\textnormal{(\Roman*)},leftmargin=13mm]
\item $A_k\ne0$. Along positive $t$, the quotient tends to $0$ when
$A_k>0$ and has unbounded modulus when $A_k<0$. Its inverse has a
logarithmic phase core and a convergent analytic lift as in
Theorem~\ref{thm:prepared}.
\item $A_k=0$ and $\beta\ne0$. The inverse has a power core, perturbed
by a convergent slope-adapted exponential transseries as in
Theorem~\ref{thm:power}.
\item $A_k=\beta=0$ and $B\ne0$. The finite limit is $C_{h,k}$.
The exact regular core is
$\tau=B^{-1}\Log(Y/C_{h,k})$, and its flat perturbation is given by
Theorem~\ref{thm:power} with $\sigma=2$.
\item $A_k=\beta=B=0$. The finite limit is $C_{h,k}$ and is infinitely
flat on closed right-half-plane subsectors. Put
$m=\ord_0(H_{h,k}-1)$ and
$c=C_{h,k}h_m$. Then all local inverse sheets are given by
Theorem~\ref{thm:flat}, with
\[
 U(Q)=\frac{H_{h,k}(Q)-1}{h_mQ^m}.
\]
\end{enumerate}
In particular, a nonconstant infinitely flat finite nonzero limit occurs
if and only if the three linear conditions
\begin{equation}\label{eq:three-conditions}
 \sum_{\delta\mid N}\frac{r_\delta\gcd(k,\delta)^2}{\delta}=0,
 \qquad \sum_{\delta\mid N}r_\delta=0,
 \qquad \sum_{\delta\mid N}\delta r_\delta=0
\end{equation}
hold.
\end{theorem}
\begin{proof}
The alternatives are exhaustive and disjoint by their definitions.
For (I), put $\epsilon=\operatorname{sign}(A_k)$ and invert
$Y=\F_{\vct r}^{\epsilon}$ instead of $\F_{\vct r}$ when necessary.
The prepared data are $A=\epsilon A_k>0$, $C=C_{h,k}^{\epsilon}$,
$\beta'=\epsilon\beta$ and
$U(t,E)=\e^{\epsilon Bt}H_{h,k}(E)^{\epsilon}$, with single action
$a>0$. The reciprocal target change is locally invertible and restores
the original target at the end. Theorem~\ref{thm:prepared} applies.
Alternatives (II) and (III) are exactly \eqref{eq:power-map}.
For (IV), Lemma~\ref{lem:H-nonconstant} gives a finite $m$ and nonzero
$h_m$; its displayed $U$ is a holomorphic unit. This proves applicability
of Theorem~\ref{thm:flat}.

For necessity of the final assertion, compare \eqref{eq:cusp-normal}
along $t\downarrow0$. A finite nonzero limit forces $A_k=0$ and
$\beta=0$, since both are real. The remaining algebraic factor is
$\e^{Bt}$. Flatness forces its linear coefficient $B$ to vanish.
Conversely, all three vanishing leaves a nonconstant exponentially small
tail. This is flat uniformly wherever $\Re(1/t)\ge c/|t|$.
\end{proof}

The original inverse in the $q$ plane is simply
$q(Y)=\zeta\exp(-t(Y))$. Analytic composition preserves the local sheet
information and the proved error estimates, up to a bounded derivative
factor on each sufficiently small chart. We give formulas in $t$ because
it is the natural cusp coordinate and displays the transseries scales.

\subsection{Minimum support and all four-scale solutions at \texorpdfstring{$q=1$}{q=1}}
At $q=1$, the three conditions are
\begin{equation}\label{eq:flat-q1}
 \sum_j r_j/\delta_j=0,\qquad
 \sum_jr_j=0,\qquad \sum_j\delta_jr_j=0.
\end{equation}

\begin{theorem}[Four-factor minimum]\label{thm:four}
On $s$ distinct positive integer scales $\delta_1,\ldots,\delta_s$, the
rational solution space of \eqref{eq:flat-q1} has dimension
$\max(s-3,0)$. In particular, a nonzero integer solution requires at
least four scales, and every set of four distinct scales admits one.
For four scales all solutions are multiples of
\begin{equation}\label{eq:kernel-four}
 r_i=K\frac{\delta_i}
              {\prod_{j\ne i}(\delta_i-\delta_j)},\qquad i=1,2,3,4.
\end{equation}
The scalar $K$ may be chosen to make all four exponents integers and
then primitive by dividing their common divisor.
\end{theorem}
\begin{proof}
Put $x_i=r_i/\delta_i$. The three rows become
$\sum x_i\delta_i^j=0$ for $j=0,1,2$.
The resulting $3\times s$ Vandermonde matrix has rank $\min(3,s)$:
every square minor of size at most three on distinct scales is nonzero.
This proves the dimension and the lower bound.
For four scales, the leading-coefficient identity in Lagrange polynomial
interpolation gives
\[
 \sum_{i=1}^4
 \frac{\delta_i^j}{\prod_{l\ne i}(\delta_i-\delta_l)}=0,
 \qquad j=0,1,2.
\]
It proves \eqref{eq:kernel-four}; the one-dimensional kernel proves
exhaustiveness. The entries are nonzero rational numbers, so clearing
denominators gives the claimed integer vector.
\end{proof}
This minimum concerns distinct Euler factors at $q=1$, not arbitrary
special-function representations or every other cusp. At other cusp
denominators the first row in \eqref{eq:three-conditions} changes and
its rank must be checked afresh.


\subsection{Intrinsic ramification and unbounded complexity with four factors}
An arbitrarily chosen common multiple $N$ can make the number $m$ in
Theorem~\ref{thm:flat} artificially large. The following criterion
separates this coordinate effect from actual ramification of the
\emph{denominator correction}.

\begin{theorem}[Exact correction ramification]\label{thm:ramification}
Let $H(Q)=1+h_mQ^m+\cdots$, $h_m\ne0$, be a nonconstant holomorphic
germ, and put
\[
 g=\gcd\{n\ge1:[Q^n](H(Q)-1)\ne0\}.
\]
Then $g\mid m$. In the flat inverse with
$u=(Y-C)/(Ch_m)$, the correction $V(u^{1/m})$ of
Theorem~\ref{thm:flat} has minimal finite ramification degree $m/g$.
The logarithmic term still has infinite monodromy and is not included
in this finite-degree assertion.

For a flat Euler quotient at $q=1$, take $N$ to be the least common
multiple of its active scales. Then $g=1$ and this degree is exactly
\begin{equation}\label{eq:intrinsic-degree}
 \frac{N}{\delta_{\max}}.
\end{equation}
Every positive integer occurs as this degree for a flat quotient with
exactly four distinct Euler factors. Four factors are minimal.
\end{theorem}
\begin{proof}
For an $m$th root of unity $\omega$, the condition
$V(\omega s)=V(s)$ is equivalent, by $V(0)=0$, to
$\psi(\omega s)=\omega\psi(s)$. Since $\psi$ is the inverse of
$f(Q)=Q((H(Q)-1)/(h_mQ^m))^{1/m}$, this is equivalent to
$f(\omega Q)=\omega f(Q)$, and hence to $H(\omega Q)=H(Q)$.
The last equivalence in the reverse direction uses the unique root
with constant term $1$. The rotation stabilizer of $H$ inside the
$m$th roots of unity is exactly the group of $g$th roots of unity:
its coefficients require $\omega^n=1$ for every occupied exponent.
Thus exactly $m/g$ circuits in $u$ are needed to restore the correction.
Equivalently its series lies in $\C\{u^{g/m}\}$ and in no smaller
finite ramification. This includes $V=0$, for which $H-1$ is a single
monomial and $g=m$.

For completeness, for any nontrivial finite Euler quotient
$H(Q)=\prod_e\Poch(Q^e)^{r_e}$ the gcd of the support of $H-1$ equals
the gcd of its active $e$. One divisibility is immediate by writing
every factor as a series in $Q^{\gcd e}$. For the converse, let $d$
divide the entire support of $H-1$. Then $\log H$ also has support in
$d\N$. If some active $e$ is not divisible by $d$, choose the smallest
such $e$. Every smaller active scale that divides $e$ would have to be
divisible by $d$, which is impossible. The coefficient of $Q^e$ in
$\log H$ is therefore $-r_e\ne0$, a contradiction.
At $q=1$ the active scales of $H$ are $e=N/\delta$. Their gcd is $1$
when $N$ is the least common multiple, and
\eqref{eq:first-q1} gives \eqref{eq:intrinsic-degree}.

For every integer $m\ge3$, choose the four scales
$1,2,m,m+1$. They are distinct, their least common multiple is
$m(m+1)$, and their maximum is $m+1$. Choose the nonzero integer kernel
vector given by Theorem~\ref{thm:four}. It has all four entries nonzero,
is flat, and has intrinsic correction degree $m$. Degree $1$ is realized
by scales $1,2,3,6$, and degree $2$ by $2,3,4,6$, again with their
nonzero kernel vectors. Minimality follows from Theorem~\ref{thm:four}.
\end{proof}
This yields unbounded branch complexity with a fixed number of product
factors. It also shows precisely which notion of ramification is meant:
finite ramification of the analytic phase correction, superimposed on
an already infinitely sheeted logarithmic inverse.

\subsection{Arithmetic action ratios in the essential regime}
There is a complementary restriction when $A_k\ne0$. Set
\[
 M_k=\sum_{\delta\mid N}r_\delta(N/\delta)\gcd(k,\delta)^2\in\Z.
\]
Then $A_k=\pi^2M_k/(6k^2N)$. After replacing the target by its reciprocal
when $M_k<0$, the ratio between the dual action and the dominant action is
\begin{equation}\label{eq:rational-action}
 \lambda=\frac{a}{|A_k|}=\frac{24}{|M_k|}\in\Q_{>0}.
\end{equation}
With $\epsilon=\operatorname{sign}(M_k)$,
$\beta'=\epsilon\beta$, $C'=C_{h,k}^{\epsilon}$, and target
$Y=\F_{\vct r}^{\epsilon}$, equation~\eqref{eq:action-transport} gives
\begin{equation}\label{eq:algebraic-Q}
 Q^{|M_k|}
 =\left(\frac{Y}{C'|A_k|^{\beta'}}\right)^{24}
        x_0^{24\beta'},\qquad 24\beta'=-12\epsilon R_0\in\Z.
\end{equation}
Thus the small dual parameter is algebraic over the pair $(Y,x_0)$;
all its fractional powers can be handled on a finite algebraic cover.
This does not eliminate the logarithmic core or assert that the full
inverse is algebraic. It does rule out irrational action-ratio
ramification for this finite integer Euler-product family, in contrast
to Proposition~\ref{prop:irrational} for general prepared transseries.

\section{Two explicit infinitely flat \texorpdfstring{$q$}{q}-analogs}
\label{sec:examples}

\subsection{A four-factor quotient with an unramified correction}
Consider
\begin{equation}\label{eq:example6}
 F_6(q)=\frac{\Poch(q^2)^5\Poch(q^6)}{\Poch(q)\Poch(q^3)^5},
 \qquad Y(t)=\frac89F_6(\e^{-t}).
\end{equation}
Its exponents on $1,2,3,6$ are $(-1,5,-5,1)$ and satisfy
\eqref{eq:flat-q1}. The cusp constant is $C_{0,1}=9/8$. With
\[
 a=\frac{2\pi^2}{3},\qquad Q=\e^{-a/t},
\]
Theorem~\ref{thm:cusp} gives the exact identity
\begin{equation}\label{eq:H6}
 Y(t)=H(Q)=\frac{\Poch(Q)\Poch(Q^3)^5}
                         {\Poch(Q^2)^5\Poch(Q^6)}.
\end{equation}
Equivalently,
$F_6(\e^{-t})F_6(\e^{-a/t})=9/8$.
This self-reciprocity is a consequence of the classical eta law, not a
new modular transformation.

The first terms, computed by \eqref{eq:logH}--\eqref{eq:H-rec}, are
\begin{align}
 H(Q)={}&1-Q+4Q^2-10Q^3+20Q^4-39Q^5\notag\\
       &+76Q^6-140Q^7+244Q^8-415Q^9+O(Q^{10}).
 \label{eq:H6-series}
\end{align}
Put $u=1-Y$. The inverse $Q=\psi(u)$ and its logarithmic unit are
\begin{align}
 \psi(u)={}&u+4u^2+22u^3+140u^4+965u^5+6992u^6\notag\\
 &+52420u^7+402880u^8+3155410u^9+25083336u^{10}+O(u^{11}),
 \label{eq:psi6}\\
 V(u)={}&4u+14u^2+\frac{220}{3}u^3+451u^4
        +\frac{15124}{5}u^5+\frac{64274}{3}u^6+O(u^7).
 \label{eq:V6}
\end{align}
All these series converge near zero. The complete inverse family is
\begin{equation}\label{eq:t6}
 \boxed{\quad t_\ell(Y)=
 -\frac{2\pi^2/3}{\Log(1-Y)+2\pi\ii\ell+V(1-Y)},
 \qquad \ell\in\Z.\quad}
\end{equation}
The real small-positive-$t$ inverse uses $u>0$ and $\ell=0$.
The numerator $2\pi^2/3$ and the full analytic denominator correction
are invisible in the algebraic expansion $Y\sim1$.

For the same quotient at the four divisor cusp denominators, the
normalized actions are
\begin{equation}\label{eq:example-actions}
 \left(\frac{6k^2}{\pi^2}A_k\right)_{k=1,2,3,6}=(0,8,-12,0).
\end{equation}
Thus the cusps of denominators $1$ and $6$ are flat, the denominator-$2$
cusp has $A_2=\pi^2/3>0$, and the denominator-$3$ cusp has
$A_3=-2\pi^2/9<0$. At $\zeta_6=\e^{\pi\ii/3}$, the constant is
$C_{1,6}=1$ and the first transformed coefficient gives
\begin{equation}\label{eq:sixth-cusp}
 F_6(\zeta_6\e^{-t})
 =1+\e^{-\pi\ii/3}\exp\left(-\frac{\pi^2}{9t}\right)
   +o\left(\exp\left(-\frac{\pi^2}{9t}\right)\right)
\end{equation}
along positive $t$. Here $a=\pi^2/54$ and the first nonzero index is
$m=6$. In fact all transformed exponents are multiples of $6$, so the
primitive coordinate is $\e^{-\pi^2/(9t)}$. This illustrates why an
auxiliary Puiseux order must not automatically be called intrinsic.

\subsection{A four-factor quotient with genuine square-root correction}
Consider the distinct scales $2,3,4,6$ and exponents $(-1,4,-4,1)$:
\begin{equation}\label{eq:example12}
 F_{12}(q)=\frac{\Poch(q^3)^4\Poch(q^6)}
                         {\Poch(q^2)\Poch(q^4)^4},
 \qquad C=\frac{16\sqrt3}{27},\qquad
 Y(t)=F_{12}(\e^{-t})/C.
\end{equation}
Here $N=12$, $a=\pi^2/3$, and
\begin{equation}\label{eq:H12}
 Y(t)=H(Q)=\frac{\Poch(Q^2)\Poch(Q^4)^4}
                        {\Poch(Q^3)^4\Poch(Q^6)},\qquad Q=\e^{-a/t}.
\end{equation}
The smallest powers are $2$ and $3$, so the primitive action coordinate
cannot remove their relative square-root correction. Specifically,
\begin{align}
 H(Q)={}&1-Q^2+4Q^3-5Q^4-4Q^5+19Q^6\notag\\
        &-20Q^7-9Q^8+60Q^9-76Q^{10}+O(Q^{11}).
 \label{eq:H12-series}
\end{align}
Put $u=1-Y$ and fix $s=\sqrt u$ coherently with
$\log s=\tfrac12\Log u$. The analytic inverse of
$s=Q\sqrt{(1-H(Q))/Q^2}$ is
\begin{align}
 \psi(s)={}&s+2s^2+\frac{15}{2}s^3+32s^4+\frac{1203}{8}s^5
             +746s^6\notag\\
 &+\frac{61467}{16}s^7+20320s^8+
       \frac{14038275}{128}s^9+601374s^{10}+O(s^{11}),
 \label{eq:psi12}\\
 V(s)={}&2s+\frac{11}{2}s^2+\frac{59}{3}s^3+\frac{337}{4}s^4
          +\frac{7843}{20}s^5+\frac{11567}{6}s^6+O(s^7).
 \label{eq:V12}
\end{align}
All inverse sheets are
\begin{equation}\label{eq:t12}
 \boxed{\quad t_\ell(Y)=-\frac{\pi^2/3}
 {\tfrac12\Log(1-Y)+\pi\ii\ell+
                  V((-1)^\ell\sqrt{1-Y})},\qquad \ell\in\Z.\quad}
\end{equation}
One circuit of $Y$ around $1$ changes the sign of the square root and
increments the sheet. Two circuits restore the square-root sign but add
$2\pi\ii$ to the logarithm. Treating the square root and logarithm as
independent principal branches would lose this compatibility.

Both examples have leading real inverse
$t\sim-2\pi^2/(3\Log(1-Y))$. Their leading logarithmic scale is therefore
identical, but one denominator correction starts at $1-Y$ and the other
at $\sqrt{1-Y}$. A leading-equivalent inverse does not determine the
subdominant ramification.

\section{Recovering the quotient from its cusp actions}
\label{sec:recovery}
An individual cusp may hide every algebraic coefficient, and some cusps
may have zero essential action. Nevertheless, the actions at the
\emph{divisor cusps together} determine the original finite exponent
vector.

For every $k\mid N$, set
\begin{equation}\label{eq:normalized-actions}
 a_k=\frac{6k^2}{\pi^2}A_k
 =\sum_{\delta\mid N}\frac{r_\delta}{\delta}\gcd(k,\delta)^2.
\end{equation}
Let $\mu$ denote the M\"obius function and
\[
 J_2(d)=\sum_{e\mid d}\mu(d/e)e^2
       =d^2\prod_{p\mid d}(1-p^{-2})>0
\]
be the second Jordan totient.

\begin{theorem}[Double M\"obius recovery]\label{thm:recover}
For every $\delta\mid N$,
\begin{equation}\label{eq:recover}
 \boxed{\quad
 r_\delta=\delta
 \sum_{\substack{m\mid N\\\delta\mid m}}
 \frac{\mu(m/\delta)}{J_2(m)}
       \sum_{e\mid m}\mu(m/e)a_e.\quad}
\end{equation}
The linear map from $(r_\delta)_{\delta\mid N}$ to $(a_k)_{k\mid N}$
is invertible, with determinant
\begin{equation}\label{eq:det}
 \prod_{d\mid N}\frac{J_2(d)}{d}
\end{equation}
when rows and columns use the same increasing divisor order.
Consequently a nontrivial Euler quotient cannot have zero essential
action at every divisor cusp, and in particular cannot be infinitely
flat at all of them.
\end{theorem}
\begin{proof}
The divisor identity $n^2=\sum_{d\mid n}J_2(d)$ gives
\begin{align*}
 a_k
 &=\sum_{d\mid k}J_2(d)
       \sum_{\substack{\delta\mid N\\d\mid\delta}}
                   \frac{r_\delta}{\delta}.
\end{align*}
M\"obius inversion on the divisors of $k$ therefore gives
\[
 b_d:=\frac1{J_2(d)}\sum_{e\mid d}\mu(d/e)a_e
     =\sum_{\substack{\delta\mid N\\d\mid\delta}}
                     \frac{r_\delta}{\delta}.
\]
Apply M\"obius inversion a second time, now to this upper-divisor sum:
$r_\delta/\delta=\sum_{\delta\mid m\mid N}\mu(m/\delta)b_m$.
This is \eqref{eq:recover}.

For the determinant, let $E_{k,d}=1$ when $d\mid k$ and $0$ otherwise.
With divisors increasing, $E$ is triangular with diagonal $1$. The
matrix in \eqref{eq:normalized-actions} factors as
\[
 E\,\operatorname{diag}(J_2(d))\,E^{\mathsf T}
          \operatorname{diag}(1/\delta).
\]
Taking determinants proves \eqref{eq:det}. If all $a_k$ vanish, recovery
gives $\vct r=0$, proving the final statements.
\end{proof}

This theorem does not assume modular-form invariance of the quotient.
It is a finite arithmetic consequence of the exact cusp actions. The
underlying divisor-matrix factorization is elementary; the contribution
here is its explicit use as an inverse-data reconstruction procedure,
not a claim to have discovered M\"obius inversion or modular rigidity.

The actions can themselves be read from leading radial inverse data.
Indeed, on a continuously chosen logarithm along the radial image,
\begin{equation}\label{eq:inverse-action-data}
 \lim_{t\downarrow0} t\log\F_{\vct r}(\zeta\e^{-t})=-A_k.
\end{equation}
Equivalently, along the corresponding inverse branch,
$t(Y)\log Y\to-A_k$. The real version
$t\log|\F_{\vct r}(\zeta\e^{-t})|\to-A_k$ avoids the logarithmic
phase convention. Even in the zero-action cases the limit exists and
is zero. Exact action data at all $k\mid N$ therefore recover every
exponent. No analogous recovery from a \emph{single} flat cusp's leading
constant is asserted. For $F_6$, substituting \eqref{eq:example-actions}
into \eqref{eq:recover} returns $(-1,5,-5,1)$.

\section{A nonmodular application: fixed-parameter shifted factorials}
\label{sec:shifted}
The exact modular tail is an advantage of the Euler-product family, not
an assumption that should silently be transferred to every $q$-analog.
We now give an all-orders inverse in a case where the algebraic series
need not converge. This makes the formal/analytic distinction concrete.

Fix $z\in(0,1)$ and define
\[
 P_z(t)=(z;\e^{-t})_\infty=\prod_{n\ge0}(1-z\e^{-nt}),\qquad \Re t>0.
\]
The following constants are positive real when appropriate:
\[
 A=\Li_2(z)>0,\quad C=\sqrt{1-z}>0,\quad
 b_j=\frac{B_{2j}}{(2j)!}\Li_{2-2j}(z).
\]
Here $B_n$ are the Bernoulli numbers with $B_2=1/6$, $B_4=-1/30$,
and $\Li_s(z)=\sum_{n\ge1}z^n/n^s$ for $|z|<1$.

\begin{proposition}[All-orders inverse of a shifted factorial]
\label{prop:shifted}
Uniformly for $z$ in a compact subinterval of $(0,1)$ and $t\to0$ in
$|\arg t|\le\pi/2-\epsilon$, for every fixed $M\ge0$,
\begin{equation}\label{eq:shifted-asymptotic}
 \log P_z(t)=-\frac{A}{t}+\log C
             -\sum_{j=1}^M b_jt^{2j-1}+O(t^{2M+1}).
\end{equation}
For $Y=P_z(t)$, choose $L=\Log(C/Y)$ on the corresponding logarithmic
cover. The local inverse has the all-orders expansion
\begin{equation}\label{eq:shifted-reversion}
 t\sim\sum_{n\ge1}\frac{L^{-n}}{n}
       [w^{n-1}]\left(A+\sum_{j\ge1}b_jw^{2j}\right)^n.
\end{equation}
Every coefficient involves only finitely many $b_j$, and only odd powers
of $L^{-1}$ occur. In particular,
\begin{align}
 t={}&\frac{A}{L}+\frac{A^2\Li_0(z)}{12L^3}\notag\\
 &+\frac1{L^5}\left(
       \frac{A^3\Li_0(z)^2}{72}-\frac{A^4\Li_{-2}(z)}{720}\right)
       +O(L^{-7}).\label{eq:shifted-first}
\end{align}
These are asymptotic assertions, not a claim of convergence of
\eqref{eq:shifted-reversion} or of a specified Borel sum.
\end{proposition}
\begin{proof}
The absolutely convergent logarithmic expansion gives
\begin{equation}\label{eq:shifted-logsum}
 \log P_z(t)=-\sum_{n\ge1}\frac{z^n}{n(1-\e^{-nt})}.
\end{equation}
On a fixed closed sector $|\arg w|\le\pi/2-\epsilon$, the kernel has
\[
 \frac1{1-\e^{-w}}
 =\frac1w+\frac12+
   \sum_{j=1}^M\frac{B_{2j}}{(2j)!}w^{2j-1}
   +R_M(w),\qquad |R_M(w)|\le C_{M,\epsilon}|w|^{2M+1}.
\]
Near zero this follows from its convergent Taylor expansion after
removing the pole. For $|w|\ge1$, the denominator is bounded away from
zero because $\Re w\ge |w|\sin\epsilon$, and the same bound follows
by bounding the finitely many subtracted terms. Inserting $w=nt$ into
\eqref{eq:shifted-logsum} bounds the remainder by
$C_{M,\epsilon}|t|^{2M+1}\sum_{n\ge1}z^nn^{2M}$, uniformly on the
stated compact sets. Summing the individual terms proves
\eqref{eq:shifted-asymptotic}.

Formally, the logarithmic equation is
$L=A/t+\sum_{j\ge1}b_jt^{2j-1}$, or
\[
 t=L^{-1}\left(A+\sum_{j\ge1}b_jt^{2j}\right).
\]
Lagrange inversion proves \eqref{eq:shifted-reversion} in the formal
power-series ring. Parity of the parenthesis proves the odd-power claim.
The first coefficients are $A$, $b_1A^2$, and
$2b_1^2A^3+b_2A^4$, which give \eqref{eq:shifted-first}.

To identify this formal inverse with an analytic asymptotic inverse,
put $x=A/t$. The exact forward logarithmic phase is
$x+O(x^{-1})$, with derivative $1+O(x^{-2})$ on smaller large-$x$
sectors. The derivative estimate follows either by differentiating the
normally convergent logarithmic sum on interior sectors or by Cauchy's
estimate applied to the uniform remainder on a slightly larger sector.
For large $L$, Lemma~\ref{lem:certificate} on a disk about $L$ of radius
$K/|L|$, with $K$ sufficiently large, gives the unique root close to $L$.
At every fixed order, substitution of the formal inverse cancels the
corresponding terms of \eqref{eq:shifted-asymptotic}; the same certificate
transports the remaining phase error into an inverse error. Taking
arbitrarily large fixed $M$ proves the full asymptotic statement.
\end{proof}

The forward expansion is part of the classical polylogarithm--Bernoulli
$q$-factorial calculus \cite{mcintosh}; we included its proof to make
the inverse statement and its sectorial remainder independent of a
formal citation. The recent gamma-product approach of Arabi Ardehali
and Rosengren provides a broader route to nonmodular asymptotics
\cite{ar}. At $z\to1$ our constants and remainder bounds are not uniform:
$C\to0$ and the negative-index polylogarithms diverge. A simultaneous
$z=\e^{-\alpha t}\to1$ limit requires a different, typically gamma-based,
core. Substituting $z=1$ into \eqref{eq:shifted-first} is invalid.

\section{Finite Gaussian multinomials: ordinary critical-point reversal}
\label{sec:finite}
For comparison with the infinitely flat boundary case, let
$n_1,\ldots,n_s$ be nonnegative integers with sum $n$, and consider the
Gaussian multinomial polynomial
\begin{equation}\label{eq:qmultinomial}
 P(q)=\begin{bmatrix}n\\n_1,\ldots,n_s\end{bmatrix}_q
 =\frac{\prod_{j=1}^n(1-q^j)}
        {\prod_{i=1}^s\prod_{j=1}^{n_i}(1-q^j)}.
\end{equation}
These are finite $q$-analogs; the standard $q$-binomial notation and
product identities are recorded in DLMF \S17.2 \cite{dlmfq}.

\begin{proposition}[Root-of-unity order and inverse sheets]
\label{prop:finite}
At a primitive $k$th root $\zeta$, the vanishing order of $P$ is exactly
\begin{equation}\label{eq:finite-order}
 \nu=\left\lfloor\frac nk\right\rfloor
       -\sum_{i=1}^s\left\lfloor\frac{n_i}{k}\right\rfloor,
 \qquad 0\le\nu\le s-1.
\end{equation}
If $\nu>0$, write
$P(\zeta+w)=c w^\nu U(w)$ with $c\ne0$ and $U(0)=1$.
There are $\nu$ ordinary Puiseux inverse sheets, obtained by evaluating
\begin{equation}\label{eq:finite-inverse}
 w(z)=\sum_{j\ge1}\frac{z^j}{j}[w^{j-1}]U(w)^{-j/\nu},
 \qquad z=(Y/c)^{1/\nu},
\end{equation}
at the $\nu$ choices of the root. The series is convergent near zero.
\end{proposition}
\begin{proof}
Each factor $1-q^j$ vanishes simply at $\zeta$ precisely when $k\mid j$.
Counting numerator and denominator factors proves the order formula;
all remaining local units have nonzero values, so no further cancellation
can change it. Write $n_i=k a_i+b_i$ with $0\le b_i<k$; then
$\nu=\lfloor\sum_i b_i/k\rfloor$, proving the bounds. For $\nu>0$,
the analytic function $wU(w)^{1/\nu}$ has derivative $1$ at zero.
Its inverse exists and Lagrange inversion gives \eqref{eq:finite-inverse}.
\end{proof}
If $\nu=0$, the polynomial has a nonzero value at the cusp, but its
derivative need not be nonzero. To study its inverse at that value one
must find the first nonzero derivative of $P(q)-P(\zeta)$; the floor
formula alone does not determine that critical order.

For example,
\begin{equation}\label{eq:finite-example}
 \begin{bmatrix}6\\2,2,2\end{bmatrix}_q
   =\Phi_3(q)^2\Phi_4(q)\Phi_5(q)\Phi_6(q),
\end{equation}
where $\Phi_j$ is the $j$th cyclotomic polynomial. At
$\zeta=\e^{2\pi\ii/3}$ this equals
$6\zeta(q-\zeta)^2+O((q-\zeta)^3)$, and the inverse branches are
\[
 q=\zeta\ \pm\sqrt{Y/(6\zeta)}+O(Y).
\]
Here the critical point has a finite-order zero and ordinary Puiseux
reversion suffices. In Section~\ref{sec:examples}, by contrast, the
boundary limit has no first nonzero algebraic derivative in the cusp
coordinate. Its logarithmic denominator is essential, not an optional
re-expression of an ordinary Taylor inverse.

\section{Sectorial transitions and exponentially small inverse jumps}
\label{sec:jumps}
Local formulas must be glued with care. Let $F_-$ and $F_+$ be
holomorphic and univalent on chosen overlapping source charts, with
local inverses $G_-$ and $G_+$. If a source transition $S$ satisfies
$F_+=F_-\circ S$, then on the corresponding target overlap
\begin{equation}\label{eq:source-transition}
 G_+=S^{-1}\circ G_-.
\end{equation}
If instead the transition acts on the target, $F_+=T\circ F_-$, then
\begin{equation}\label{eq:target-transition}
 G_+=G_-\circ T^{-1}.
\end{equation}
These exact identities follow by composition and are not interchangeable.
They apply to actual sectorial sums when the necessary summability and
univalence hypotheses have been established. They do not themselves
prove resurgence or construct a Stokes automorphism.

\begin{proposition}[All-order perturbation transport]\label{prop:jumps}
Let $F$ be univalent with inverse $G$, and let $\Delta$ be holomorphic
on a source neighborhood. For sufficiently small $\epsilon$, the inverse
of $F_\epsilon=F+\epsilon\Delta$ satisfies locally
\begin{equation}\label{eq:jump-all}
 G_\epsilon(y)=G(y)+\sum_{n\ge1}\frac{(-\epsilon)^n}{n!}
 \left(\frac{\dd}{\dd y}\right)^{n-1}
 \left[G'(y)\,\Delta(G(y))^n\right].
\end{equation}
The expansion converges on sufficiently small compact subcharts for
small $\epsilon$. In particular,
\begin{equation}\label{eq:jump-first}
 \left.\frac{\partial G_\epsilon(y)}{\partial\epsilon}\right|_{0}
       =-\frac{\Delta(G(y))}{F'(G(y))}.
\end{equation}
\end{proposition}
\begin{proof}
Put $u=F(z)$ and $d(u)=\Delta(G(u))$. Then
$y=u+\epsilon d(u)$ and $z=G(u)$. Lagrange inversion with the observable
$G$ gives \eqref{eq:jump-all}. The analytic implicit-function theorem
and a sufficiently small $\epsilon$ disk justify convergence uniformly
on compact subcharts. The $n=1$ term is \eqref{eq:jump-first}; it also
follows by differentiating $F_\epsilon(G_\epsilon(y))=y$.
\end{proof}
Sauzin's nonlinear resurgent analysis supplies rigorous settings in
which nonlinear substitutions and inversions preserve appropriate
resurgent and summability classes \cite{sauzin}. Applying those results
to a particular divergent transseries requires checking its precise
continuability and growth hypotheses. In the exact modular examples,
$H$ and its inverse germ already converge: no unspecified Borel sum is
required to justify their local formulas. Analytic monodromy of their
logarithms must still be kept distinct from changes between lateral sums
of a divergent expansion.

\subsection{An inverse action budget at a flat limit}
The first-order jump law displays a useful threshold. Suppose, along a
closed decay sector,
\[
 F(t)=Y_*+c\e^{-A_0/t}(1+O(\e^{-a/t})),\qquad
 \Delta(t)=b\e^{-B_0/t},
\]
where $A_0,a,B_0>0$ and $b,c\ne0$. Then
\begin{equation}\label{eq:jump-action-budget}
 \left.\partial_\epsilon G_\epsilon(F(t))\right|_0
 \sim-\frac{b\,t^2}{cA_0}\e^{-(B_0-A_0)/t}.
\end{equation}
This follows by differentiating the leading forward term and applying
\eqref{eq:jump-first}. Forward action $B_0$ becomes inverse action
$B_0-A_0$, dressed by $t^2$. It is therefore incorrect to transfer the
forward error exponent unchanged through an infinitely flat inverse.

When $B_0>A_0$, a bounded perturbation coefficient has an exponentially
small ratio to the leading displacement; a small inverse-phase chart
remains appropriate. When $B_0=A_0$, it changes the leading amplitude
$c$ to $c+\epsilon b$ (unless that coefficient cancels, in which case a
new leading action is needed). When $0<B_0<A_0$, any fixed nonzero
$\epsilon b$ eventually changes the leading action itself. The
linear-response formula is still a derivative at $\epsilon=0$, but is
not a uniform approximation for fixed $\epsilon$ in that last regime.
These distinctions specify an inverse error budget without conflating
a small absolute forward error with a small relative error in the
quantity that is being inverted.

\section{Algorithms, verification, and reproducibility}
\label{sec:verification}

\subsection{A branch-aware workflow}
The preceding proofs suggest a practical algorithm rather than only
existence assertions. Given an Euler quotient, a cusp, a target, and a
chosen continuation sheet, first compute the integer and cyclotomic data
in \eqref{eq:cusp-integers}--\eqref{eq:H-data}. Evaluate the three
invariants exactly and select the chart in
Theorem~\ref{thm:classification}. Next generate the required coefficients
using \eqref{eq:logH} and \eqref{eq:H-rec}. In the flat chart find the
first action, remove any redundant primitive-coordinate factor, and
compute the analytic inverse $\psi$ and $V$ with
\eqref{eq:psi-coeff}. The logarithm and root are continued together, as
in \eqref{eq:flat-inverse}; a principal logarithm is not a universal
sheet-selection algorithm.

Finally evaluate a truncated inverse and certify its chart whenever
explicit analytic bounds are available. Corollary~\ref{cor:flat-tail}
provides a closed tail estimate from $R,M$; in more general charts
Lemma~\ref{lem:certificate} converts a computable residual and a
local derivative bound into existence and error control. The existence
of analytic radii in a theorem is not by itself a numerical interval
certificate for an implementation. Deriving sharp effective radii is a
separate, useful task.

For exact coefficient computation, the following recurrence avoids
symbolic differentiation. If $A(z)=\sum a_nz^n$, $a_0=1$, and
$B(z)=A(z)^p=\sum b_nz^n$, then $b_0=1$ and
\begin{equation}\label{eq:power-rec}
 nb_n=\sum_{k=1}^n((p+1)k-n)a_kb_{n-k}.
\end{equation}
It follows by comparing coefficients in $AB'=pA'B$.
Rational $p$ and rational input coefficients therefore remain exact.
The supplied program uses rational arithmetic for this recurrence,
Lagrange coefficients, compositions, and cusp recovery.

\subsection{Numerical product tails}
For $|q|=r<1$, truncating the logarithm of $\Poch(q)$ after factor $M$
has the absolute tail bound
\begin{equation}\label{eq:product-tail}
 \left|\sum_{j>M}\log(1-q^j)\right|
 \le\frac{r^{M+1}}{(1-r)(1-r^{M+1})}.
\end{equation}
Indeed, $|\log(1-q^j)|\le\sum_{l\ge1}r^{jl}/l
\le r^j/(1-r^j)$, and summing the geometric bound proves the formula.
For quotients, multiply the individual bounds by the absolute values of
their integer exponents and add. The program chooses a cutoff using
this type of bound but evaluates the finite sum in ordinary arbitrary
precision, not with outward-rounded arithmetic. Its recorded numerical
agreement is therefore a diagnostic, not a machine-certified enclosure.

\subsection{Executed checks}
The package contains \texttt{code/verify.py}, which has no network
dependency and uses Python's rational arithmetic and \texttt{mpmath}.
The recorded run checks exact reversion compositions for both examples,
all three flatness moments, the four-node kernel formula, and recovery
on every basis vector for levels $1,2,6,12,24,30,60$. It also checks the
arbitrary-degree four-factor construction, the displayed shifted-factorial
reversion coefficients on exact rational test inputs, and the finite
Gaussian example. The recorded run passed 619 exact assertions; software versions and
individual coefficient data are recorded in \texttt{data/verification.json}.

Fifteen independent forward-versus-dual product tests use three exponent
vectors, including one with nonzero weight and essential action, at
cusps $(h,k)=(0,1),(1,2),(1,3),(2,5),(1,6)$. All use
$t=0.37+0.09\ii$ at 115 decimal digits, and all relative discrepancies
are below $10^{-90}$. This tests, in particular, the Dedekind phases in
\eqref{eq:single-cusp}, not only the positive-real $q\to1$ case.

Table~\ref{tab:inverse} evaluates the first example from its
\emph{original} product at $q=\e^{-t}$ and reconstructs $t$ from the
inverse polynomial $\psi_N(u)$, using
$-a/(\log\psi_N(u)+2\pi\ii\ell)$. Here $N$ is the degree of $\psi_N$,
not of $V_N$. Truncating $\psi$ at degree $N$ generally gives a
logarithmic-unit error $O(u^N)$, so the indexing differs by one from
Corollary~\ref{cor:flat-tail}.
\begin{table}[htbp]
\centering\small
\begin{tabular}{@{}lrrrr@{}}
\toprule
Reference $t$ & Sheet & $N=1$ error & $N=4$ error & $N=8$ error\\
\midrule
$0.5$ & $0$ & $2.93\cdot10^{-7}$ & $5.06\cdot10^{-22}$ & $2.28\cdot10^{-41}$\\
$0.35$ & $0$ & $5.10\cdot10^{-10}$ & $3.95\cdot10^{-32}$ & $2.84\cdot10^{-61}$\\
$0.25$ & $0$ & $1.41\cdot10^{-13}$ & $1.74\cdot10^{-45}$ & $1.08\cdot10^{-87}$\\
$0.4+0.08\ii$ & $1$ & $1.37\cdot10^{-8}$ & $8.15\cdot10^{-27}$ & $8.90\cdot10^{-51}$\\
\bottomrule
\end{tabular}
\caption{Absolute inverse errors from the executed non-interval checks.
The sheet is selected to match the known reference point, not inferred
from a principal-logarithm convention.}
\label{tab:inverse}
\end{table}
The complex row specifically requires a nonprincipal logarithmic sheet.
Using the principal logarithm would converge to a different inverse,
not to the displayed reference point. The test computes the matching
integer from the known reference phase. An actual inverse application
must supply its sheet by initial data or continuation rather than use
an unavailable reference solution.

Two additional checks use $s=\pm10^{-5}$ in the second example. With
$\psi$ through degree ten, the forward residuals in $1-H(\psi(s))-s^2$
are below $7\cdot10^{-54}$. These test the two distinct square-root
branches in the convergent dual germ. Neither these tests nor the exact
finite computations constitute a proof of an untested infinite theorem;
the conventional proofs above provide that part of the argument.

\section{Further research questions and topics}
\label{sec:research}
The following program distinguishes extensions with concrete technical
obstacles from claims already proved here.
\begin{enumerate}[label=\textbf{R\arabic*.},leftmargin=14mm]
\item \textbf{Countable actions and accumulation.}
Find useful complete topological algebras in which the parameter-lift
argument survives countably many actions with finite positive-projection
budgets. Determine what replaces this budget when actions accumulate at
a positive value or at zero. The formal contraction lemma is available,
but admissibility of substitution and analytic realization are not
automatic.

\item \textbf{Divergent analytic units and summation-compatible inversion.}
Replace the convergent unit $U$ in Theorem~\ref{thm:prepared} by a
specified resurgent or multisummable unit. Prove an inverse theorem with
explicit lateral-sum compatibility, and compute which singularities
appear in the Borel plane after the dominant-phase change. Established
closure theorems are starting points, not substitutes for their
hypothesis checks in each $q$-family.

\item \textbf{Confluence of essential and flat charts.}
For continuously weighted products or shifted parameters, let the
essential action tend to zero simultaneously with $t$. Construct a
uniform chart that connects alternatives (I), (II), (III), and (IV) of
Theorem~\ref{thm:classification}. A fixed nonzero-action inverse is
singular in this transition, and merely setting $A_k=0$ in its Lambert
core cannot be justified.

\item \textbf{Effective first-action bounds at arbitrary cusps.}
Lemma~\ref{lem:H-nonconstant} proves termination of a coefficient search,
but not a sharp complexity bound. Bound the first nonzero index and the
gcd of the occupied support in terms of $N,k$ and the active exponent
vector, including cancellations among cyclotomic phases. Identify a
closed form analogous to \eqref{eq:first-q1} when one exists.

\item \textbf{Minimum support away from $q=1$.}
Classify the rank drops of the three rows in
\eqref{eq:three-conditions} at arbitrary $k$. Determine exactly which
cusps admit two- or three-factor flat quotients, and classify their
intrinsic correction ramification. The four-factor result here is
specific to $q=1$.

\item \textbf{Effective inverse radii and rigorous enclosures.}
For concrete $H$, locate the nearest critical values of
$Q((H(Q)-1)/(h_mQ^m))^{1/m}$ and derive numerical $R,M$ for
Corollary~\ref{cor:flat-tail}. Build outward-rounded complex interval
certificates for both sheet continuation and inverse truncation. The
existing numerical code is a starting test suite, not such a certificate.

\item \textbf{Optimal truncation with intact action blocks.}
When algebraic blocks diverge but the action expansion converges, choose
joint algebraic and action cutoffs that meet a prescribed inverse error.
The action subtraction in \eqref{eq:jump-action-budget} shows why a
forward optimal cutoff can be suboptimal or insufficient after inversion.
A useful theorem should distinguish fixed-action accuracy from total
absolute accuracy.

\item \textbf{Moving shifted parameters and gamma cores.}
Study $(z;q)_\infty$ with $z=\e^{-\alpha t}$, including exceptional
$\alpha$ near poles and zeros of the limiting gamma factors. Combine the
recent product representation \cite{ar} with a uniform inverse chart
that remains valid as $z\to1$. Include root-of-unity analogues and
coalescing gamma singularities.

\item \textbf{Growing cusp denominator.}
The results here keep $k,N$ fixed. Determine a useful joint limit in
which $k\to\infty$ while $t\to0$ and the dual parameter
$\exp(-4\pi^2/(k^2Nt))$ remains small. Uniform control must also track
the phases, the first occupied dual index, and inverse critical values;
density of rational cusps prevents a naive global patching argument.

\item \textbf{Stable recovery from noisy cusp data.}
Analyze the condition numbers of the divisor-action matrix in
Theorem~\ref{thm:recover}. Exploit its prime-power tensor structure and
integer exponent constraints to reconstruct quotients from approximate
inverse measurements. Determine how many selected cusps suffice when
the active support is sparse or only bounded a priori.

\item \textbf{Ramification spectra with bounded exponent height.}
Four factors suffice for every correction degree, but the integer
exponents supplied by the Vandermonde kernel may be large. Determine
sharp tradeoffs between ramification degree, maximum scale, coefficient
height, and number of factors. Classify symmetries that force correction
coefficients to vanish without changing the leading action.

\item \textbf{Basic hypergeometric functions and multiple saddles.}
Extend the cusp-to-inverse strategy beyond products to basic
hypergeometric series and Nahm-type sums. Multiple competing saddles
introduce interference and possible derivative cancellation; root-of-unity
asymptotics such as those in \cite{gz} provide forward input, but inverse
branch selection must be proved independently.

\item \textbf{Global transition data.}
Combine logarithmic monodromy, finite correction ramification, and
sectorial Stokes maps in a single explicitly based system of source and
target charts. Determine when the local transitions in
\eqref{eq:source-transition} satisfy a global compatibility law. Do not
identify a dominance wall with a Stokes discontinuity without specifying
an analytic summation problem.

\item \textbf{Formalization in stages.}
A practical Lean program could first formalize the finite divisor-matrix
recovery, Vandermonde kernel, and rational coefficient recurrences; then
the complete filtered fixed-point lemma and finite power-series
identities; and finally the analytic inverse chart, error certificate,
and modular bridge. The repository's existing exact modules and their
stated limitations should be respected. No Lean formalization of the
new theorems is claimed in this package.
\end{enumerate}

\section{Conclusion}
A complex transseries inverse is not only a list of coefficients. It
comprises a dominant phase, an admissible analytic or formal algebra,
a collection of small actions, and coherent continuation data. For
prepared exponential maps these ingredients lead to a convergent
multivariable inverse. When the dominant essential action vanishes,
slope-adapted and logarithmic--Puiseux charts provide the missing
alternatives.

The $q$-product application makes this principle explicit. At every
rational cusp, a finite Euler quotient has an exact elementary factor
times a convergent dual product. Three computable invariants determine
its inverse regime. In the infinitely flat regime, an ordinary analytic
inverse in the dual variable becomes a convergent correction inside an
infinitely sheeted logarithmic denominator. Four Euler factors are both
necessary and sufficient at $q=1$, and already realize every finite
correction ramification degree. Despite local flatness, divisor-cusp
actions collectively recover the full exponent vector.

These conclusions answer the scoped questions posed in
Section~\ref{sec:scope}, with proofs and reproducible checks. They do
not establish a universal inversion theory for all transseries, a new
modular transformation law, or worldwide priority. The remaining program
is to extend the same explicit separation of algebra, asymptotics,
and analytic continuation to divergent blocks, moving parameters,
and multi-saddle $q$-analogs.

\clearpage
\appendix
\section{Proof dependencies and research-status ledger}
\label{app:ledger}
\begin{longtable}{@{}>{\raggedright\arraybackslash}p{0.23\textwidth}>{\raggedright\arraybackslash}p{0.34\textwidth}>{\raggedright\arraybackslash}p{0.37\textwidth}@{}}
\toprule
Result group & Main proof ingredients & Scope of the claim\\
\midrule
\endhead
Prepared inverse & Analytic implicit-function theorem; filtered
contraction; Lagrange inversion & Holomorphic parameter unit and finite
admissible action set; no general divergent-unit assertion.\\[4pt]
Flat inverse atlas & Ordinary analytic inversion in $Q$; logarithms on
a cover & Exact local sheets and convergent correction with explicit
Taylor-tail bound.\\[4pt]
Cusp preparation & Classical eta transformation; finite products & Exact
identity on $\Re t>0$ for integer Euler quotients; not a new modular law.\\[4pt]
Four regimes & Exact normal form; leading-term comparison & Necessary
and sufficient flatness test at each fixed cusp.\\[4pt]
Four-factor and ramification results & Vandermonde kernel; rotational
stabilizers; Euler-product coefficient uniqueness & Minimum factor count
at $q=1$; exact correction degree; every degree attained with four factors.\\[4pt]
Cusp-data recovery & Jordan divisor identity; two M\"obius inversions &
Exact finite-dimensional recovery with specified $N$ and divisor-cusp
measurements.\\[4pt]
Shifted factorial & Kernel expansion; formal reversal; complex residual
certificate & All finite asymptotic orders for fixed $z$ in compact
subsets of $(0,1)$; no $z\to1$ uniformity claim.\\[4pt]
Finite multinomial & Count simple roots; analytic Puiseux inversion &
Finite-order zeros at roots of unity; nonzero-value criticality treated
separately.\\[4pt]
Sectorial jump transport & Composition; Lagrange inversion with an
observable & Exact chart identities and perturbation series; resurgent
summation compatibility requires additional hypotheses.\\
\bottomrule
\end{longtable}
The article supplies conventional proofs for these scoped statements.
The established tools are attributed in the bibliography. The new
formulations and combined classifications are proposed research results
of this draft, not independently refereed or proof-assistant-checked
claims. The repository comparison is restricted to the cited inspected
documentation. The program's exact finite checks and floating-point
diagnostics have the distinct evidential roles explained in
Section~\ref{sec:verification}.

\clearpage
\begin{thebibliography}{99}
\raggedright
\bibitem{repo}
V. Reshetnikov, \emph{ProveIt: Transseries documentation}, repository
snapshot \texttt{8e9cd6f00e0ac79c4425ff6abdfa64b497cf3f41},
inspected 4 October 2026. Targeted documentation:
\path{Analysis/Transseries/README.md} and
\path{Analysis/Transseries/docs/series-and-transseries/README.md}.
\url{https://github.com/VladimirReshetnikov/ProveIt}.

\bibitem{repoq}
ProveIt research package, \emph{Certified Inversion Through the
$q\to1$ Transition: Signed remainders, optimal truncation, and a modular
cancellation window}, README and editorial amendments, at the snapshot
in \cite{repo}. Directory:
\path{Analysis/Transseries/docs/series-and-transseries/Certified_Inversion_q_to_1_Transition/}.

\bibitem{repohahn}
ProveIt research package, \emph{Beyond Finite-Action Folds: Critical Hahn
Transseries, Stable Sector Asymptotics, and Sharp Action Budgets}, README
and editorial amendments, at the snapshot in \cite{repo}. Directory:
\path{Analysis/Transseries/docs/series-and-transseries/Critical_Hahn_Transseries_Beyond_Finite_Action_Folds/}.

\bibitem{sauzin}
D. Sauzin, \emph{Nonlinear analysis with resurgent functions},
Annales scientifiques de l'\'Ecole Normale Sup\'erieure (4)
\textbf{48} (2015), no. 3, 667--702.
\url{https://doi.org/10.24033/asens.2255}.
Preprint: \url{https://arxiv.org/abs/1212.4477}.

\bibitem{dlmf}
NIST Digital Library of Mathematical Functions,
\emph{\S27.14: Unrestricted Partitions}, especially the eta product and
modular transformation formulas (27.14.12)--(27.14.14).
\url{https://dlmf.nist.gov/27.14}. Accessed 4 October 2026.

\bibitem{dlmfq}
NIST Digital Library of Mathematical Functions,
\emph{\S17.2: $q$-Calculus}.
\url{https://dlmf.nist.gov/17.2}. Accessed 4 October 2026.

\bibitem{mcintosh}
R. J. McIntosh, \emph{Some asymptotic formulae for $q$-shifted factorials},
The Ramanujan Journal \textbf{3} (1999), 205--214.
\url{https://doi.org/10.1023/A:1006949508631}.

\bibitem{gz}
S. Garoufalidis and D. Zagier, with an appendix by S. Zwegers,
\emph{Asymptotics of Nahm sums at roots of unity}, The Ramanujan Journal
\textbf{55} (2021), 219--238.
\url{https://doi.org/10.1007/s11139-020-00266-x}.
Preprint: \url{https://arxiv.org/abs/1812.07690}.

\bibitem{fr}
V. Fantini and C. Rella, \emph{Modular resurgence, $q$-Pochhammer
symbols, and quantum operators from mirror curves}, Letters in
Mathematical Physics \textbf{116} (2026), no. 2, article 39.
\url{https://doi.org/10.1007/s11005-026-02063-x}.
Preprint: arXiv:2506.08265 (2025; revised 1 April 2026).
\url{https://arxiv.org/abs/2506.08265}.

\bibitem{ar}
A. Arabi Ardehali and H. Rosengren, \emph{A New Product Formula for
$(z;q)_\infty$, with Applications to Asymptotics}, Constructive
Approximation (2026), version of record published 19 September 2026.
\url{https://doi.org/10.1007/s00365-026-09783-2}.
Preprint: \url{https://arxiv.org/abs/2602.11329}.
\end{thebibliography}
\end{document}
